%قبل از شروع مسئله ترجیحاً چیزی ننویسید. اما اگر خواستید بنویسید، در قالب لیستی (فقرات و شمارش) بنویسید.
%برای اضافه کردن نمره، بعدی از شروع‌مسئله، [x] را اضافه کنید و به جای x  نمره مورد نظر را قرار دهید. برای زیرمسئله‌ها هم به همین صورت است.
%\شروع‌مسئله[10]
\شروع‌مسئله
روابط بازگشتی زیر را تحلیل کنید.
\شروع{زیرمسائل}
\زیرمسئله
$T(n) = \sqrt n \times T(\sqrt n) + n$

\زیرمسئله

$T(n) = T(\frac{n}{2} + \sqrt n) + \sqrt {6046}$
\زیرمسئله
$T(n) = mT(\frac{n}{m})+ O(n)$ \footnote{با استفاده از روش استقرا این قسمت را تحلیل کنید.}

($m > 1$ و $n$ توانی از $m$)

\زیرمسئله
$T(n) = T( \ceil{ \frac{n}{2}}) + T( \floor{ \frac{n}{2}}) + O(n)$ \footnote{$n$ .لزوماً توانی از ۲ نیست}

\پایان{زیرمسائل}

\solution{
قسمت ۱:
تعریف می کنیم
$f(n) = \frac{T(n)}{n}$
حال با جایگزینی در صورت اصلی رابطه داریم
$f(n) = f(\sqrt n) + 1$
حال به استقرا اثبات می کنیم که
$\forall k : f(n) = f(n^{\frac{1}{2^k}}) + k$
به وضوح حکم برای ۱ برقرار است
حال اگر حکم برای ${1, 2, .., k}$ برقرار باشد
$f(f(n^{\frac{1}{2^k}}) = f(n^{\frac{1}{2^{k+1}}}) + 1$
$f(n) = f(n^{\frac{1}{2^k}}) + k$
$\Rightarrow f(n) = f(n^{\frac{1}{2^{k+1}}}) + k + 1$
و حکم استقرا ثابت می شود. حال اگر در صورت تعمیم یافته رابطه قرار دهیم 
$k = log(log(n))$
داریم
$f(n) = f(2) + log(log(n))$
حال با توجه به نحوه تعریف داریم
$T(n) = \theta (\sqrt n \times log(log(n)) $

 قسمت ۲:
داریم
$\exists N \forall n > N : \frac{n}{2} < \frac{n}{2} + \sqrt n < \frac{2n}{3}$
و همچنین با استفاده از استفاده از استقرا می توان ثابت کرد که $T$ صعودی است
حال داریم
$T(n) < T(\frac{2n}{3}) + \sqrt {6046})) < ... < T(\frac{2^kn}{3^k}) + k \times \sqrt{6046}$
$\Rightarrow T(n) < c_1 \times log(n)$
و از طرفی دیگر
$T(n) > T(\frac{n}{2}) + \sqrt {6046})) > ... > T(\frac{n}{2^k}) + k \times \sqrt{6046}$
$\Rightarrow T(n) > c_2 \times log(n)$
$\Rightarrow T(n) = \theta (log(n))$

قسمت ۳:
 ابتدا جملات کوچک تر را بررسی می‌کنیم تا به یک الگوی کلی دست پیدا کنیم.

$n = m \rightarrow T(m) = mT(1) + O(m) = mO(1) + O(m) = O(m) + O(m) = 2O(m)$

$n = m^{2} \rightarrow T(m^{2}) = mT(m) + O(m^{2}) = 2mO(m) = O(m^{2}) = 2O(m^{2}) + O(m^{2}) = 3O(m^{2})$

پس متوجه میشویم $T(m^{k})$ از الگوی زیر پیروی می	کند.

$T(m^{k}) = (k+1)O(m^{k})$

حال از استقرا استفاده می‌کنیم.
$T(n) = mT(\frac{n}{m})+ O(n)$


$T(m^{k}) = mT(m^{k-1}) + O(m^{k}) = mkO(m^{k-1}) + O(m^{k}) = kO(m^{k}) + O(m^{k}) = (k+1)O(m^{k})$

حال رابطه ی نهایی را بر حسب n می‌نویسیم.

$T(n) = (log^{n}_{m} + 1)O(n) = O(nlogn)$


قسمت ۴:
 رابطه ی بازگشتی نوشته شده در این قسمت همان merge sort است.
برای اثبات این قسمت هم از استقرا استفاده می‌کنیم.

$T(n) \leq T( \ceil{ \frac{n}{2}}) + T( \floor{\frac{n}{2}}) + c_{2}n$

$T(1) = c_{1}$

$T(n) \leq c_{3}nlog(n)$


در ابتدا تعریف می‌کنیم که $c_{3} = 2c_{2} + c_{1}$

پایه استقرا:


$T(2) \leq T(1) + T(1) + 2c_{2} \leq c_{1} +c_{1} + 2c_{2} \leq c_{1} + c_{3} \leq 2c_{3} \leq c_{3}2log(2)$

فرض استقرا 

$T(n) \leq T( \floor{ \frac{n}{2}}) + T( \ceil{\frac{n}{2}}) + c_{2}n \leq c_{3} \floor{ \frac{n}{2}}log\floor{ \frac{n}{2}} c_{3} \ceil{ \frac{n}{2}}log\ceil{ \frac{n}{2}} +c_{2}n \leq c_{3} \floor{ \frac{n}{2}}log\ceil{ \frac{n}{2}} c_{3} \ceil{ \frac{n}{2}}log\ceil{ \frac{n}{2}} +c_{2}n \leq c_{3}( \floor{ \frac{n}{2}} + \ceil{ \frac{n}{2}}) log\ceil{ \frac{n}{2}} + c_{2}n \leq c_{3}nlog\ceil{ \frac{n}{2}} + c_{2}n \leq c_{3}nlog\ceil{ \frac{n}{2}} + \frac{c_{3}}{2}n \leq c_{3}n(log\ceil{ \frac{n}{2}} + \frac{1}{2}) \leq c_{3}n(log\ceil{ \frac{n}{2}} \sqrt{2}) \leq c_{3}nlog(n)$




}

